\section{Driving Logic}
\subsubsection*{Things to Write}

\begin{itemize}
  \item If you are targeting both challenges, you must include separate discussions of how your driving logic matches each course.
  \item Design of the lane following and obstacle avoidance systems must be specifically described, along with how the vehicle navigates through its environment. Describe how the system uses GPS for waypoint navigation and localization.
  \item If targeting AutoNav, describe how your robot uses the perception system’s output and other information to follow lanes, detect and avoid obstacles, climb the ramp, and navigate to the GPS waypoints. Include how vehicle motion is generated and monitored to move the vehicle through the course, including how it deals with complex obstacles including switchbacks and center islands dead ends, traps, and potholes.
  \item If targeting SelfDrive, describe how the vehicle uses the perception system’s output and otherinformation to follow the rules of the road while navigating the course that includes pedestrians,stop signs, road obstacles, and other items.
  \item Clearly identify how your design was driven by the driving logic requirements you listed.
  \item For each driving logic requirement you listed, compare the target value with the most recentlymeasured actual value, and discuss what this tells you about your vehicle
\end{itemize}

\subsection*{Path Planning}

The car’s path planning logic is dependent on perception outputs as well as the behavior required by each challenge. The current methodology focuses on local path waypoint generation using an A* planner operating on a Nav2 costmap that is generated from and will be dynamically updated by LiDAR data.

The perception pipeline publishes a Nav2 occupancy costmap representing nearby obstacles detected by the LiDAR sensor. Occupied territories are marked as high cost cells which correspond to obstacles like barrels, cones, potholes, and so on, while free space cells represent parts of the road that the car is allowed to travel on. This costmap serves as the primary input to the local planner. While the current version of the costmap contains only obstacle clusters from the LiDAR, in the future, it will be integrated with the outputs generated by lane segmentation, marking the boundaries of the current lane as an object, as well as using the centerline of the current lane as part of the goal for A* to target.

The local planner uses the A* search algorithm to compute a path through the environment. The algorithm first finds a cell near the center of the local costmap which represents the car’s current position that doesn’t contain any obstacles and marks this as the starting position. Then, to find a goal destination for A* to generate a path for, the planner searches for cells some target distance ahead in the forward direction and chooses a destination cell without obstacles that meets this criteria. This method is used rather than having a globally fixed coordinate since we’re unaware of the obstacles placed on the track in some of the driving challenges, and the obstacles might change each lap, so searching for a goal that is obstacle free in the general direction of movement is preferable.

In the future, this goal selection process can be modified to also consider the state of the finite state machine that controls the behavior tree of the car to choose an appropriate destination. For example, when turns are necessary, the forward direction destination selection process would not be compatible, so we would need to consider a different method that can choose a destination in the direction of the turn.

The A* algorithm, to find a path, looks at neighboring cells using a cost function determined by both traversal distance and occupancy cost. Cells with obstacle costs above a threshold are treated as non-traversable, while lower cost cells, such as the inflation layer of the costmap, are still available for the car to travel on but are penalized according to their occupancy value. This encourages the planner to maintain safe clearance from nearby obstacles instead of driving directly along obstacle boundaries.

Once a valid path is generated from the start cell to the destination cell, the path of grid points is converted into world frame coordinates using the Nav2 costmap metadata. To extract waypoints from the A* path, a spline is fitted through the path, since raw A* output can contain sharp corners and uneven motion. The spline is generated using SciPy interpolation tools and produces a set of waypoints that can be used to control the car. The spline output represents the car’s local path trajectory and is published as a ROS topic. These waypoints generated by the spline will then be inputted into the MPC, which will generate the steering and throttle commands necessary to ensure the car can physically travel in the intended path.

The path planning logic will also, in the future, be combined with the finite state machine controlling the car, represented as a Nav2 behavior tree, which will ensure that behavior dependent on the YOLO object identification system, such as ensuring the car halts when a stop sign is identified, is taken into consideration when generating waypoints.
